\section{Power Series Solution}\label{power_series_solution}
\subsection{Differential Equation}\label{diff_eq_harmonic_oscillator}

The simple harmonic oscillator is an interesting system that allows us to model many complex systems. Classically, the total energy $E$ of a particle of mass $m$ moving in one dimension with a restoring force $F=-kx$ is given by $$E=\frac{p^2}{2m}+\frac{1}{2}m\omega^2x^2,$$where $\omega=\sqrt{\frac{k}{m}}$ is the angular frequency of the oscillations.

\begin{framedrmk*}[Ubiquitousness of the harmonic oscillator]
    Simple harmonic oscillations occur in quadratic potentials. However, \textit{any} potential is approximated as such to first order around a minimum. For example, consider an arbitrary $V(x)$ with a minimum at $x_0$. For $x\approx x_0$, the Taylor expansion of $V(x)$ gives $$V(x)=V(x_0)+(x-x_0)V'(x_0)+\frac{1}{2}(x-x_0)^2V''(x_0)+\mathcal{O}((x-x_0)^3).$$Since $x_0$ is a minimum, $V'(x_0)=0$. Thus, the potential is approximately quadratic for $x\approx x_0$, and small-amplitude oscillations about this point can be approximated with a harmonic oscillator.
\end{framedrmk*}

To define an analogous \textit{quantum} harmonic oscillator, we define the Hamiltonian as $$\hat{H}=\frac{\hat{p}^2}{2m}+\frac{1}{2}m\omega^2\hat{x}^2.$$The harmonic oscillator potential is $$V(x)=\frac{1}{2}m\omega^2x^2.$$We wish to find the energy eigenstates, which will all be bound states since the potential diverges in both directions. There are two methods to solving this system. The first, which we will discuss presently, is a power series solution to the Schrödinger equation. It has the virtue of being applicable to a vast range of potentials. The second, which will be covered in Part \ref{algebraic_solution}, is an ingenious algebraic solution that contains more intuition.

We begin with the analytical method by writing the Schrödinger equation for a state $\phi$: $$-\frac{\hbar^2}{2m}\frac{d^2\phi}{dx^2}+\frac{1}{2}m\omega^2x^2\phi=E\phi.$$As a first step, we change all the constants to be dimensionless. We begin by introducing a unit-free coordinate $u$ to replace the position $x$. We set $x=au$, where $a$ is a constant with units of length. By dimensional analysis, the only combination of $\hbar$, $m$, and $\omega$ with the correct dimensions is $$a=\sqrt{\frac{\hbar}{m\omega}}.$$Substituting $x=au$ into the Schrödinger equation yields $$-\frac{\hbar^2}{2ma^2}\frac{d^2\phi}{du^2}+\frac{1}{2}m\omega^2a^2u^2\phi=E\phi.$$Here, we note that since $\phi$ appears in every term, its units cancel out. Furthermore, since $u$ is dimensionless, we must have $$[\frac{\hbar^2}{ma^2}]=[m\omega^2a^2]=[E].$$Indeed, substituting the value of $a$, we see that both of these quantities equal $\hbar\omega$, which correctly has units of energy. Thus, we have $$-\frac{1}{2}\hbar\omega\frac{d^2\phi}{du^2}+\frac{1}{2}\hbar\omega u^2\phi=E\phi.$$We can even cancel the remaining units by multiplying both sides by $\frac{2}{\hbar\omega}$, which gives $$-\frac{d^2\phi}{du^2}+u^2\phi=\mathcal{E}\phi,$$where we defined the unit-free energy $\mathcal{E}=\frac{E}{\hbar\omega}$. Rearranging, we arrive at a very nice dimensionless form of the Schrödinger equation: $$\boxed{\frac{d^2\phi}{du^2}=(u^2-\mathcal{E})\phi}.$$

\subsection{Recursion Relation}\label{recursion_relation_harmonic_oscillator}

This differential equation surely has solutions for all values of $\mathcal{E}$. Quantization arises because these solutions are not normalizable except for specific values of $\mathcal{E}$. To understand this intuitively, we examine the equation in the limit as $u$ goes to infinity. Then we have $$\lim_{u\to\infty}\phi''(u)= u^2\phi(u).$$Since we want a $u^2$ factor to appear through twice differentiation, we can try an exponential in $u^2$: $$\phi(u)\sim u^ke^{\alpha u^2/2}.$$Then to leading term in $u$, we have $\phi'(u)\approx\alpha u\phi$, and $$\lim_{u\to\infty}\phi''=(\alpha u)^2\phi.$$from which we conclude $\alpha=\pm 1$. Thus, our general solutions for very large $u$ are $$\phi\approx Au^ke^{-u^2/2}+Bu^ke^{u^2/2}.$$The second term diverges as $u\rightarrow\pm\infty$, so we only want the first term. Also, note that the $u^k$ term played no role in the analysis, so we can replace it with an arbitrary polynomial $h(u)$: $$\phi(u)=h(u)e^{-u^2/2}.$$Effectively, we've captured the tendency of the solutions to the Schrödinger equation to diverge in the $e^{-u^2/2}$ term. 

Substituting $\phi(u)$ into the differential equation, we obtain yet another differential equation, this time in $h$: $$\frac{d^2h}{du^2}-2u\frac{dh}{du}+(\mathcal{E}-1)h=0.$$

We can express $h(u)$ as its Taylor series expansion: $$h(u)=\sum_{k=0}^\infty a_ku^k.$$We consider the $u^j$ terms in the differential equation and declare that their coefficient must vanish. The double derivative produces a $(j+2)(j+1)a_{j+2}u^j$ term. Similarly, the second term produces a $-2ja_j$ component, and the last term produces a component of $(\mathcal{E}-1)a_j$. Since the $u^j$ terms must sum to $0$, we have $$(j+2)(j+1)a_{j+2}=(2j+1-\mathcal{E})a_j.$$This is a two-step recursion relation for the Taylor series coefficients $a_k$: $$\boxed{a_{j+2}=\frac{2j+1-\mathcal{E}}{(j+2)(j+1)}a_j}.$$Since this relation skips a step, the value of $a_0$ uniquely defines the even coefficients and thus the even solutions, and $a_1$ defines the odd solutions.
\newpage
\begin{framedrmk*}[Two independent variables]
    Note that $a_0=h(u)$ and $a_1=h'(u)$. It makes sense that we need both $a_0$ and $a_1$ to define $h(u)$, since a specific solution to a second-order differential equation is specified by both the value and derivative of the function at a point.
\end{framedrmk*}

\subsection{Energy Levels and Eigenstates}\label{energy_levels_eigenstates_harmonic_oscillator}

We now show that the Taylor series expansion of $h(u)$ must terminate. Consider the ratio of $a_{j+2}$ to $a_j$ for very large $j$: $$\lim_{j\to\infty}\frac{a_{j+2}}{a_j}\approx\frac{2}{j}.$$Note that $$e^{u^2}=\sum_{n=0}^\infty\frac{1}{n!}u^{2n}=\sum_{j\text{ even}}\frac{1}{(j/2)!}u^j.$$This series has coefficients $c_j=\frac{1}{(j/2)!}$ for even $j$, and we see that as $j$ goes to $\infty$, we have $$\frac{c_{j+2}}{c_j}=\frac{(j/2)!}{(j+2)/2!}=\frac{2}{j+2}\approx\frac{2}{j}.$$ Thus, if $h(u)$ does not terminate, the wavefunction $\phi$ behaves like $$\phi(u)=h(u)e^{-u^2/2}\sim e^{u^2}e^{-u^2/2}= e^{u^2/2},$$which is the diverging solution from before. We conclude that $h(u)$ has some finite degree $j$.

However, recall that we identified a recursion relation for the $a_j$ (boxed on the previous page). For this relation to give $a_j=0$, we must choose $\mathcal{E}$ such that $$2j+1-\mathcal{E}=0.$$This equation quantizes the energy as $$\mathcal{E}=2n+1,\quad n\in\mathbb{N}_0.$$We have $$h(u)=a_nu^n+a_{n-2}u^{n-2}+\ldots$$Since we only have bound states for the harmonic oscillator, Corollary \ref{even_potential_cor} tells us that all solutions are even or odd. If $n$ is odd, we have an odd solution. If $n$ is even, we have an even solution. 

The energy levels are given by $$\boxed{E_n=\hbar\omega(n+\frac{1}{2}).}$$Note that for the harmonic oscillator, the energy levels are evenly spaced with common differences $\hbar\omega$. The corresponding polynomials $h_n(u)$ are called the \textbf{Hermite polynomials}. 

\begin{frameddefn}[Hermite polynomials]
    The Hermite polynomial $H_n(u)$ has leading term $2^nu^n$ by convention and satisfies the differential equation $$\frac{d^2}{du^2}H_n-2u\frac{d}{du}H_n+2nH_n.$$The first few Hermite polynomials are \begin{align*}H_0(u)&=1,\\H_1(u)&=2u,\\H_2(u)&=4u^2-2,\\H_3(u)&=8u^3-12u.\end{align*}The generating function for the Hermite polynomials is an exponential with formal parameter $z$: $$e^{-z^2+2zu}=\sum_{u=0}^\infty\frac{z^n}{n!}H_n(u).$$
\end{frameddefn}

For completeness, we write the energy eigenstates in terms of $x$. Recall that $u^2=\frac{x^2}{a^2}$, where $a^2=\frac{\hbar}{m\omega}$. Thus, we have $$\boxed{\phi_n(x)=N_nH_n\bigg(x\sqrt{\frac{m\omega}{\hbar}}\bigg)e^{-\frac{m\omega}{2\hbar}x^2},}$$where $N_n$ is a normalization constant.

\section{Algebraic Solution}\label{algebraic_solution}
\subsection{Factoring the Hamiltonian}\label{factoring_hamiltonian}

In Part \ref{power_series_solution}, we showed a general analytical solution to the quantum harmonic oscillator that can be applied to a wide variety of potentials. Here, we demonstrate a more elegant approach to the problem. Specifically, we begin by trying to \textit{factor} the Hamiltonian, which we rewrite below: $$\hat{H}=\frac{\hat{p}^2}{2m}+\frac{1}{2}m\omega^2\hat{x}^2=\frac{1}{2}m\omega^2(\hat{x}^2+\frac{\hat{p}^2}{m^2\omega^2}).$$What do we mean by ``factor''? We want to find an operator $\hat{V}$ such that $$H=V^\dagger V+\#.$$Note that $V^\dagger V$ is Hermitian, since $(V^\dagger V)^\dagger=V^\dagger V^{\dagger\dagger}=V^\dagger V$.  The arbitrary constant at the end does not change the properties of the Hamiltonian.

Note that $a^2+b^2=(a-ib)(a+ib)$. Therefore, we consider \begin{align*}(\hat{x}-\frac{i\hat{p}}{m\omega})(\hat{x}+\frac{i\hat{p}}{m\omega})&=\hat{x}^2+\frac{\hat{p}^2}{m^2\omega^2}+\frac{i}{m\omega}(\hat{x}\hat{p}-\hat{p}\hat{x})\\&=\hat{x}^2+\frac{\hat{p}^2}{m^2\omega^2}+\frac{i}{m\omega}[\hat{x},\hat{p}]\\&=\hat{x}^2+\frac{\hat{p}^2}{m^2\omega^2}-\frac{\hbar}{m\omega}.\end{align*}We define $\hat{V}=\hat{x}+\frac{i\hat{p}}{m\omega}$, and since $\hat{x}$ and $\hat{p}$ are Hermitian, we have $\hat{V}^\dagger=\hat{x}-\frac{i\hat{p}}{m\omega}$. We can thus write \begin{align*}\hat{H}&=\frac{1}{2}m\omega^2(V^\dagger V+\frac{\hbar}{m\omega})\\&=\frac{1}{2}m\omega^2V^\dagger V+\frac{1}{2}\hbar\omega.\end{align*}Also note that the commutator of $V$ and $V^\dagger$ is simple: \begin{align*}[V,V^\dagger]&=[\hat{x}+\frac{i\hat{p}}{m\omega},\hat{x}-\frac{i\hat{p}}{m\omega}]=-\frac{i}{m\omega}[\hat{x},\hat{p}]+\frac{i}{m\omega}[\hat{p},\hat{x}].\end{align*}Substituting $[\hat{x},\hat{p}]=i\hbar$, we have $$[V,V^\dagger]=\frac{2\hbar}{m\omega}.$$This implies that $$\bigg[\sqrt{\frac{m\omega}{2\hbar}}V,\sqrt{\frac{m\omega}{2\hbar}}V^\dagger\bigg]=1.$$We now define the unit-free operators $$\hat{a}=\sqrt{\frac{m\omega}{2\hbar}}V,\quad\hat{a}^\dagger=\sqrt{\frac{m\omega}{2\hbar}}V^\dagger.$$For reasons that will become clear later, we call $\hat{a}$ the \textbf{annihilation operator} and $\hat{a}^\dagger$ the \textbf{creation operator}. From the definitions of $V$ and $V^\dagger$ above, we can express $\hat{a}$ and $\hat{a}^\dagger$ in terms of $\hat{x}$ and $\hat{p}$, 
$$\boxed{\hat{a}=\sqrt{\frac{m\omega}{2\hbar}}(\hat{x}+\frac{i\hat{p}}{m\omega}),\quad\hat{a}^\dagger=\sqrt{\frac{m\omega}{2\hbar}}(\hat{x}-\frac{i\hat{p}}{m\omega}),}$$and vice versa:$$\boxed{\hat{x}=\sqrt{\frac{\hbar}{2m\omega}}(\hat{a}+\hat{a}^\dagger),\quad\hat{p}=i\sqrt{\frac{m\omega\hbar}{2}}(\hat{a}^\dagger-\hat{a}).}$$

We now have $$V^\dagger V=\frac{2\hbar}{m\omega}\hat{a}^\dagger\hat{a}.$$Substituting into the Hamiltonian, we get $$\boxed{\hat{H}=\hbar\omega(\hat{a}^\dagger\hat{a}+\frac{1}{2}).}$$We call $\hat{a}^\dagger\hat{a}$ the \textbf{number operator} $\hat{N}$. Since it effectively encodes for the Hamiltonian, it's probably a good idea to figure out how $\hat{N}$ interacts with other operators. We can quickly see that \begin{align*} [\hat{N},\hat{a}]=[\hat{a}^\dagger\hat{a},\hat{a}]=[\hat{a}^\dagger,\hat{a}]\hat{a}=-\hat{a},\\ [\hat{N},\hat{a}^\dagger]=[\hat{a}^\dagger\hat{a},\hat{a}^\dagger]=\hat{a}^\dagger[\hat{a},\hat{a}^\dagger]=\hat{a}^\dagger. \end{align*}Notice that the annihilation operator $\hat{a}$ is accompanied by a negative sign, while the creation operator $\hat{a}^\dagger$ has a positive sign. By induction, we can show that \begin{align*}[\hat{N},(\hat{a})^k]=-k(\hat{a})^k,\\ [\hat{N},(\hat{a}^\dagger)^k]=k(\hat{a}^\dagger)^k.\end{align*}This suggests the reason why $\hat{N}$ is called the number operator: its commutator with some number of creation or annihilation operators is the same operator multiplied by ($\pm$) the number of operators. It almost seems like an eigenvalue equation, where $k$ would be the eigenvalue.

Finally, we note the commutation relations \begin{align*}[\hat{a}^\dagger,(\hat{a})^k]=-k(\hat{a})^{k-1},\\ [\hat{a},(\hat{a}^\dagger)^k]=k(\hat{a}^\dagger)^{k-1}.\end{align*}

\begin{framedrmk*}[Supersymmetric quantum mechanics]
    The process of factorizing the Hamiltonian can be applied to many problems. In fact, there are entire books describing these types of problems, which are grouped together under the field of \textbf{supersymmetric quantum mechanics}.
\end{framedrmk*}

\subsection{Ground State Wavefunction}\label{ground_state_harmonic}

Consider some arbitrary normalized wavefunction $\psi$. We compute the expectation value of $\hat{H}$ on $\psi$: $$\langle\hat{H}\rangle_\psi=(\psi,\hat{H}\psi)=(\psi,\hbar\omega\hat{a}^\dagger\hat{a}\psi+\frac{1}{2}\hbar\omega\psi).$$Splitting the sum, $$\langle\hat{H}\rangle_\psi=\hbar\omega(\psi,\hat{a}^\dagger\hat{a}\psi)+\frac{1}{2}\hbar\omega(\psi,\psi).$$Since $\psi$ is normalized, $(\psi,\psi)=1$. Also, by the definition of the Hermitian conjugate, $(\psi,\hat{a}^\dagger\hat{a}\psi)=(\hat{a}\psi,\hat{a}\psi)$. Thus, we have $$\langle\hat{H}\rangle_\psi=\hbar\omega(\hat{a}\psi,\hat{a}\psi)+\frac{1}{2}\hbar\omega\geq\frac{1}{2}\hbar\omega.$$Here we see the great benefit of the factorized Hamiltonian: since the inner product of any function with itself is nonnegative, we can immediately conclude that every energy eigenstate has energy at least $\frac{1}{2}\hbar\omega$! The minimum energy will be realized for the ground state $\psi_0$ if $(\hat{a}\psi_0,\hat{a}\psi_0)=0$, which implies $$\hat{a}\psi_0=0.$$In other words, the operator $\hat{a}$ \textit{annihilates} the ground state, and is thus called the annihilation operator. From the definition of $\hat{a}$, $$\bigg(x+\frac{i}{m\omega}\frac{\hbar}{i}\frac{d}{dx}\bigg)\psi_0(x)=\bigg(x+\frac{\hbar}{m\omega}\frac{d}{dx}\bigg)\psi_0(x)=0.$$Note that through this process, we have turned a second-order differential equation into a first-order one. Rearranging, $$\frac{d\psi_0}{dx}=-\frac{m\omega}{\hbar}x\psi_0.$$Solving, $$\boxed{\psi_0(x)=\bigg(\frac{m\omega}{\pi\hbar}\bigg)^{\frac{1}{4}}e^{-\frac{m\omega}{2\hbar}x^2}.}$$We can check that $\psi_0$ is an energy eigenstate with energy $E_0=\frac{1}{2}\hbar\omega$: $$\hat{H}\psi_0=\hbar\omega(\hat{a}^\dagger\hat{a}+\frac{1}{2})\psi_0=\frac{1}{2}\hbar\omega\psi_0.$$

\subsection{Excited State Wavefunctions} \label{excited_states_harmonic_oscillator}

We now define the state $\psi_1=\hat{a}^\dagger\psi_0$. This state is an energy eigenstate if it is a number eigenstate, so we consider $$\hat{N}\psi_1=\hat{N}\hat{a}^\dagger\psi_0=[\hat{N},\hat{a}^\dagger]\psi_0,$$because $\hat{a}^\dagger\hat{N}\psi_0=\hat{a}^\dagger\hat{a}^\dagger\hat{a}\psi_0=0$. We have a formula for this commutator: $$\hat{N}\psi_1=[\hat{N},\hat{a}^\dagger]\psi_0=\hat{a}^\dagger\psi_0=\psi_1.$$Thus, $\psi_1$ is a number eigenstate with eigenvalue $1$. It is thus also an energy eigenstate with eigenvalue $E=\hbar\omega(N+\frac{1}{2})=\frac{3}{2}\hbar\omega$.  

We also ask if $\psi_1$ is normalized. We see that$$(\psi_1,\psi_1)=(\hat{a}^\dagger\psi_0,\hat{a}^\dagger\psi_0)=(\psi_0,\hat{a}\hat{a}^\dagger\psi_0)=(\psi_0,[\hat{a},\hat{a}^\dagger]\psi_0)=(\psi_0,\psi_0)=1.$$Thus, $\psi_1$ is automatically normalized as long as $\psi_0$ is. Next consider the state $\psi_2'=\hat{a}^\dagger\hat{a}^\dagger\psi_0=(\hat{a}^\dagger)^2\psi_0$. Again, we operate $\hat{N}$ on $\psi_2'$: $$\hat{N}\psi_2'=\hat{N}(\hat{a}^\dagger)^2\psi_0=[\hat{N},(\hat{a}^\dagger)^2]\psi_0=2(\hat{a}^\dagger)^2=2\psi_2'.$$Thus, $\psi_2'$ is a number eigenstate and also an energy eigenstate, with energy eigenvalue $\frac{5}{2}\hbar\omega$. Again, we try normalizing $\psi_2'$: \begin{align*}(\psi_2'\psi_2') &= (\hat{a}^\dagger\hat{a}^\dagger\,\psi_0,\hat{a}^\dagger\hat{a}^\dagger\,\psi_0)\\ &=(\psi_0,\hat{a}\hat{a}\hat{a}^\dagger\hat{a}^\dagger\,\psi_0) \\ &=(\psi_0,\hat{a}\,[\hat{a},(\hat{a}^\dagger)^2]\,\psi_0) \\ &=(\psi_0,2\hat{a}\hat{a}^\dagger\,\psi_0) \\ &=2(\psi_0,[\hat{a},\hat{a}^\dagger]\,\psi_0)=2.\end{align*}The normalization constant is $\frac{1}{\sqrt{2}}$, so the proper definition of the second excited state is $$\psi_2=\frac{1}{\sqrt{2}}\hat{a}^\dagger\hat{a}^\dagger\psi_0.$$In general, it turns out that $$\boxed{\psi_n=\frac{1}{\sqrt{n!}}(\hat{a}^\dagger)^n\psi_0.}$$The operator $\hat{a}^\dagger$ creates excited states out of the ground state, and is thus called the creation operator. Since they are energy eigenstates with different eigenvalues, the $\psi_n$ are orthonormal: $(\psi_m,\psi_n)=\delta_{mn}$.

Applying the number operator, $$\hat{N}\psi_n=\frac{1}{\sqrt{n!}}\hat{N}(\hat{a}^\dagger)^n\psi_0=\frac{1}{\sqrt{n!}}[\hat{N},(\hat{a}^\dagger)^n]\psi_0=\frac{1}{\sqrt{n!}}n(\hat{a}^\dagger)^n\psi_0=n\psi_n.$$Thus, the $n$th excited state has energy $E=\hbar\omega(n+\frac{1}{2})$. We can also check normalization: $$(\psi_n,\psi_n)=\frac{1}{n!}((\hat{a}^\dagger)^n\psi_0,(\hat{a}^\dagger)^n\psi_0)=\frac{1}{n!}(\psi_0,(\hat{a})^n(\hat{a}^\dagger)^n\psi_0).$$Now consider replacing the last $\hat{a}(\hat{a}^\dagger)^n$ by the commutator $[\hat{a},(\hat{a}^\dagger)^n]=n(\hat{a}^\dagger)^{n-1}$, which we can do because $(\hat{a}^\dagger)^n\hat{a}\psi_0=0$. This brings out a factor of $n$ and gets rid of one factor of $\hat{a}$ and one factor of $\hat{a}^\dagger$. Thus, repeating this process gives us an overall factor of $n!$, making $\psi_n$ normalized.

\subsection{Creation and Annihilation Operators on Energy Eigenstates}\label{creation_annihilation_operators}

We now formalize the effect of $\hat{a}$ and $\hat{a}^\dagger$ on the energy eigenstates. Consider $$\hat{a}\psi_n=\frac{1}{\sqrt{n!}}\hat{a}(\hat{a}^\dagger)^n\psi_0=\frac{1}{\sqrt{n!}}[\hat{a},(\hat{a}^\dagger)^n]\psi_0=\frac{1}{\sqrt{n!}}n(\hat{a}^\dagger)^{n-1}\psi_0.$$By definition, $(\hat{a}^\dagger)^{n-1}\psi_0=\sqrt{(n-1)!}\psi_{n-1}$, so $$\hat{a}\psi_n=\frac{n}{\sqrt{n!}}\sqrt{(n-1)!}\,\psi_{n-1}=\sqrt{n}\,\psi_{n-1}.$$Similarly, we have $$\hat{a}^\dagger\psi_n=\frac{1}{\sqrt{n!}}(\hat{a}^\dagger)^{n+1}\psi_0=\frac{1}{\sqrt{n!}}\sqrt{(n+1)!}\,\psi_{n+1}=\sqrt{n+1}\,\psi_{n+1}.$$Collecting the results, we have We now formalize the effect of $\hat{a}$ and $\hat{a}^\dagger$ on the energy eigenstates. Consider $$\hat{a}\psi_n=\frac{1}{\sqrt{n!}}\hat{a}(\hat{a}^\dagger)^n\psi_0=\frac{1}{\sqrt{n!}}[\hat{a},(\hat{a}^\dagger)^n]\psi_0=\frac{1}{\sqrt{n!}}n(\hat{a}^\dagger)^{n-1}\psi_0.$$By definition, $(\hat{a}^\dagger)^{n-1}\psi_0=\sqrt{(n-1)!}\psi_{n-1}$, so $$\hat{a}\psi_n=\frac{n}{\sqrt{n!}}\sqrt{(n-1)!}\,\psi_{n-1}=\sqrt{n}\,\psi_{n-1}.$$Similarly, we have $$\hat{a}^\dagger\psi_n=\frac{1}{\sqrt{n!}}(\hat{a}^\dagger)^{n+1}\psi_0=\frac{1}{\sqrt{n!}}\sqrt{(n+1)!}\,\psi_{n+1}=\sqrt{n+1}\,\psi_{n+1}.$$Collecting the results, we have
\begin{align*}
\hat{a}\psi_n&=\sqrt{n}\,\psi_{n-1},\\ \hat{a}^\dagger\psi_n &= \sqrt{n+1}\,\psi_{n+1}.
\end{align*}
These relations make it clear that the annihilation operator decreases the number of an energy eigenstate, while the creation operator increases the number.

\begin{framedexpl}[Expectation values in a harmonic oscillator]
    Find the expectation values of $\hat{x}$ and $\hat{p}$ in the state $\psi_n$.
\end{framedexpl}

\begin{framedansw*}
    Consider the value $\langle\hat{x}\rangle_{\psi_n}=\int x\psi_n(x)^2 dx$. Since the potential is even, $\psi_n$ is either even or odd. Either way, $\psi_n^2$ is even. However, $x$ is odd, so the integral cancels out and we have $\langle\hat{x}\rangle_{\psi_n}=0$.

    It's also worth it to consider this expectation value using the creation and annihilation operators. By definition, we have $\hat{x}=\sqrt{\frac{\hbar}{2m\omega}}(\hat{a}+\hat{a}^\dagger)$. Thus, we have $$\langle\hat{x}\rangle=(\psi_n,\hat{x}\psi_n)=\sqrt{\frac{\hbar}{2m\omega}}(\psi_n,(\hat{a}+\hat{a}^\dagger)\psi_n).$$As we found above, the second term is $\sqrt{n}\,\psi_{n-1}+\sqrt{n+1}\,\psi_{n+1}$. Thus, $\langle\hat{x}\rangle\sim(\psi_n,\psi_{n-1})+(\psi_n,\psi_{n+1})=0$.

    Furthermore, note that $\psi_n$ is a stationary state. Thus, the expectation value of $x$ must be time-independent, so there can be no expected momentum. We conclude that $\langle\hat{p}\rangle_{\psi_n}=0$ also.
\end{framedansw*}

\begin{framedexpl}[Uncertainty in a harmonic oscillator]
    Find the position uncertainty $\Delta x$ in $\psi_n$.
\end{framedexpl}

\begin{framedansw*}
    We have $(\Delta x)^2=\langle\hat{x}^2\rangle-\langle\hat{x}\rangle^2=\langle\hat{x}^2\rangle$, since the second term goes to $0$.

    Expanding, we see that \begin{align*}\langle\hat{x}^2\rangle&=(\psi_n,(\hat{x})^2\psi_n)\\ &=\frac{\hbar}{2m\omega}(\psi_n,(\hat{a}+\hat{a}^\dagger)^2\psi_n)\\ &=\frac{\hbar}{2m\omega}(\psi_n,(\hat{a}\hat{a}^\dagger+\hat{a}^\dagger\hat{a}^\dagger+\hat{a}\hat{a}^\dagger+\hat{a}^\dagger\hat{a})\psi_n).\end{align*}Note that the $\hat{a}\hat{a}$ and $\hat{a}^\dagger\hat{a}^\dagger$ terms result in the wrong number eigenstate, so the inner product of these states with $\psi_n$ is $0$. This leaves $$\langle\hat{x}^2\rangle=\frac{\hbar}{2m\omega}(\psi_n,(\hat{a}\hat{a}^\dagger+\hat{a}^\dagger\hat{a})\psi_n]).$$Furthermore, note that $\hat{a}^\dagger\hat{a}=\hat{N}$. The other order, $\hat{a}\hat{a}^\dagger$, can be written as $[\hat{a},\hat{a}^\dagger]+\hat{a}^\dagger\hat{a}=1+\hat{N}$. Therefore, we have $$\langle\hat{x}^2\rangle=\frac{\hbar}{2m\omega}(\psi_n,(1+2\hat{N})\psi_n)=\frac{\hbar}{2m\omega}(\psi_n,\psi_n)(1+2n).$$Our final result is $$\boxed{(\Delta x)^2=\langle\hat{x}^2\rangle_{\psi_n}=\frac{\hbar}{m\omega}(n+\frac{1}{2}).}$$
\end{framedansw*}